% Solution: Tadelis (2013) Exercise 11.1 "Disagreement". Input by \solitem.
% Rewritten from ../pratica-lecture-06-solutions/sol-tadelis-11-1.tex, deepened. Check: check_F_L06.py.
\begin{exsol}{Tadelis Exercise 11.1}{236}{Foundations}{1}
\asked (i) A Nash equilibrium of the ultimatum game ($T = 1$) whose outcome is no agreement, with
payoffs $(0,0)$; (ii) the same outcome as a Nash equilibrium for any number $T$ of periods.

\method Nash equilibrium asks for best replies \emph{given the strategies}, so a responder may
threaten to reject offers that are never made, and even the offer that is made if it leaves her
indifferent. Subgame perfection then removes such threats. Tadelis \S11.1 (ultimatum, pp.~220--224)
and \S11.2 (finitely many rounds, pp.~224--227). Slides: Lecture~6, frames~13--14 (the ultimatum
game, strategies $S_1 = [0,1]$, $S_2 = \{s_2:[0,1]\to\{A,R\}\}$, and the backward-induction
argument) and frames~15--18 (alternating offers); the SPE of every finite game is unique
(frame~19). Lecture~4, frame on non-credible threats, for the Nash/SPE contrast.

\solution
\textbf{Notation.} $x$ is player~1's share and $1 - x$ player~2's (as in frame~13). A strategy of
player~2 in the ultimatum game is a whole response function $s_2 : [0,1] \to \{A, R\}$.

\textbf{(i) $T = 1$.} Take
\[
  s_1^* : \ x = 1, \qquad s_2^*(x) = R \ \text{ for every } x \in [0,1].
\]
Outcome: player~1 claims the whole pie, player~2 rejects, payoffs $(0,0)$.
\begin{itemize}
  \item \emph{Player~1.} Every offer is rejected under $s_2^*$, so every $x$ gives him $0$. Offering
  $x = 1$ is a best reply (all offers tie).
  \item \emph{Player~2.} Only her response to the offer actually made, $x = 1$, affects her payoff:
  accepting gives $1 - 1 = 0$, rejecting gives $0$. Rejecting is a best reply (a tie). Her responses
  to other offers are never reached and change nothing.
\end{itemize}
No player gains by deviating, so $(s_1^*, s_2^*)$ is a Nash equilibrium with payoffs $(0,0)$.

Note the role of $x = 1$: if player~1 offered $x < 1$, rejecting would cost player~2 $1 - x > 0$ and
``reject everything'' would no longer be a best reply on the path.

\textbf{(ii) Any $T$.} In every period $t = 1, \dots, T$ let the proposer claim the whole pie for
himself and let the responder reject every offer. Under this profile agreement is never reached and
both get $0$. Check unilateral deviations:
\begin{itemize}
  \item a proposer who changes his offer is still rejected (the responder rejects everything), so he
  still gets $0$;
  \item a responder who accepts instead of rejecting gets the share the opponent leaves him, which is
  $0$, discounted: $\delta^{t-1}\cdot 0 = 0$;
  \item any combination of such changes by one player leaves him at $0$, because his opponent's
  strategy never gives him a positive share.
\end{itemize}
So the profile is a Nash equilibrium for every $T$ and every $\delta \in (0,1)$, with payoffs $(0,0)$.

\textbf{Why it is not subgame perfect.} Consider the last period (period~$T$, or the only period
when $T = 1$) after an offer that leaves the responder $1 - x > 0$. In that subgame accepting gives
$\delta^{T-1}(1-x) > 0$ and rejecting gives $0$, so ``reject'' is not a best reply: the threat is not
credible. Backward induction instead gives immediate agreement (frames~16--19).

\interp Nash equilibrium alone does not discipline bargaining: by threatening to reject, either side
can support disagreement or (by similar constructions) any division of the pie. That is why
the lecture solves bargaining games by subgame perfection: it pins down a unique outcome with
immediate agreement, and the responder's acceptance rule must be optimal after every offer, not just
after the one on the path.

\verify \texttt{check\_F\_L06.py} computes, for $T = 1,\dots,6$ and $\delta \in \{\frac12,
\frac9{10}\}$, the best payoff a unilateral deviator can get against ``claim everything / reject
everything'' (on a grid of offers and acceptance rules) and finds $0$; it also checks that against a
claim $x < 1$ rejecting is strictly worse for the responder. Output: ``Tadelis 11.1: disagreement is
Nash for T = 1..6 (best deviation payoff 0); x < 1 would break it''. Tadelis prints no answer.

\mistakes
\begin{itemize}
  \item Calling the profile subgame perfect because nobody deviates on the path. Subgame perfection
  also tests the responder after offers that are never made, and there ``reject $1 - x > 0$'' fails.
  \item Letting player~1 offer a positive share to player~2. Then rejecting costs her $1 - x > 0$, so
  it is not even a Nash equilibrium. The offer must leave her indifferent ($x = 1$).
  \item Describing player~2's strategy as ``reject'' (one action). Her strategy is a response to
  \emph{every} offer $x \in [0,1]$; the equilibrium needs ``reject all $x$''.
  \item For $T > 1$, forgetting the later proposers: if a later proposer offered a positive share,
  an earlier responder might want to wait for it, and the construction must make every proposer claim
  everything.
\end{itemize}

\answer{$T = 1$: player~1 offers $x = 1$ (keeps everything) and player~2 rejects every offer; no
one can gain, so it is a Nash equilibrium with outcome ``no agreement'', payoffs $(0,0)$. For any
$T$: every proposer claims the whole pie and every responder rejects every offer; a deviator still
gets $0$. These equilibria are not subgame perfect: in the last period rejecting a positive share is
not a best reply.}
\end{exsol}
